\section{2013 Analysis \& Differential Equations}

\subsection{Individual}

\begin{exercise}
\label{analysis:2013:1}
Suppose that $f$ is an integrable function on $\mathbb{R}^d$. For each $\alpha>0$, let $E_{\alpha}=\{x\mid |f(x)|>\alpha\}$. Prove that:
\[
\int_{\mathbb{R}^d}|f(x)|\,dx = \int_0^{\infty} m(E_{\alpha})\,d\alpha.
\]
\end{exercise}

\begin{solution}
\textbf{Prerequisites:}
\begin{itemize}
    \item Measure Theory: Definition of Lebesgue integral.
    \item Tonelli's Theorem: Swapping the order of integration for non-negative functions.
    \item Characteristic functions $\chi_E(x)$.
\end{itemize}

\textbf{Steps:}
\begin{enumerate}
    \item We express the absolute value $|f(x)|$ as an integral of its characteristic function over height $\alpha$:
    \[
    |f(x)| = \int_0^{|f(x)|} 1 \, d\alpha = \int_0^\infty \chi_{\{ \alpha < |f(x)| \}}(\alpha) \, d\alpha.
    \]
    \item Substitute this expression into the original integral over $\mathbb{R}^d$:
    \[
    \int_{\mathbb{R}^d} |f(x)| \, dx = \int_{\mathbb{R}^d} \left( \int_0^\infty \chi_{\{ \alpha < |f(x)| \}}(\alpha) \, d\alpha \right) dx.
    \]
    \item Since the integrand is non-negative and measurable, we apply Tonelli's Theorem to swap the order of integration:
    \[
    \int_{\mathbb{R}^d} |f(x)| \, dx = \int_0^\infty \left( \int_{\mathbb{R}^d} \chi_{\{ x \mid |f(x)| > \alpha \}}(x) \, dx \right) d\alpha.
    \]
    \item The inner integral is simply the measure of the set $E_\alpha = \{x \mid |f(x)| > \alpha\}$:
    \[
    \int_{\mathbb{R}^d} \chi_{E_\alpha}(x) \, dx = m(E_\alpha).
    \]
    \item Substituting this back, we obtain the desired identity:
    \[
    \int_{\mathbb{R}^d} |f(x)| \, dx = \int_0^\infty m(E_\alpha) \, d\alpha.
    \]
\end{enumerate}

\textbf{Intuition:}
Think of the integral as the "volume" under the graph of $|f(x)|$. We can compute this volume by adding up horizontal slices. At each height $\alpha$, the "width" of the graph is the measure of the set where $|f(x)|$ exceeds $\alpha$. Summing these widths over all heights gives the total volume.
\end{solution}

\begin{exercise}
\label{analysis:2013:2}
Let $p(z)$ be a polynomial of degree $d\geq 2$, with distinct roots $a_1,a_2,\cdots,a_d$. Show that
\[
\sum_{i=1}^d \frac{1}{p'(a_i)} = 0.
\]
\end{exercise}

\begin{solution}
\textbf{Prerequisites:}
\begin{itemize}
    \item Complex Analysis: Residue Theorem.
    \item Rational functions and their behavior at infinity.
    \item Residue formula for simple poles: $\text{Res}(f, a) = \frac{1}{p'(a)}$ for $f(z) = \frac{1}{p(z)}$.
\end{itemize}

\textbf{Steps:}
\begin{enumerate}
    \item Define the function $f(z) = \frac{1}{p(z)}$. Since the roots $a_i$ of $p(z)$ are distinct, $f(z)$ has simple poles at each $z = a_i$.

    \item Calculate the residue at each pole:
    \[
    \text{Res}(f, a_i) = \lim_{z \to a_i} (z - a_i) \frac{1}{p(z)} = \frac{1}{p'(a_i)}.
    \]

    \item Let $C_R$ be a large circle of radius $R$ enclosing all the roots. By the Residue Theorem:
    \[
    \oint_{C_R} \frac{1}{p(z)} \, dz = 2\pi i \sum_{i=1}^d \text{Res}(f, a_i) = 2\pi i \sum_{i=1}^d \frac{1}{p'(a_i)}.
    \]

    \item Analyze the integral as $R \to \infty$. Since $p(z)$ is a polynomial of degree $d$, $|p(z)| \approx C|z|^d$ for large $|z|$. Thus:
    \[
    \left| \oint_{C_R} \frac{1}{p(z)} \, dz \right| \leq (2\pi R) \cdot \frac{1}{C R^d} = \frac{2\pi}{C R^{d-1}}.
    \]

    \item Since $d \geq 2$, $d-1 \geq 1$, so the integral tends to 0 as $R \to \infty$:
    \[
    \lim_{R \to \infty} \oint_{C_R} \frac{1}{p(z)} \, dz = 0.
    \]

    \item Therefore, $2\pi i \sum_{i=1}^d \frac{1}{p'(a_i)} = 0$, which implies $\sum_{i=1}^d \frac{1}{p'(a_i)} = 0$.
\end{enumerate}

\textbf{Intuition:}
The sum of the residues of a rational function in the complex plane is related to its residue at infinity. For a function that decays sufficiently fast ($1/z^2$ or faster), the sum of all its finite residues must be zero because there is no residue "leakage" at infinity.
\end{solution}

\begin{exercise}
\label{analysis:2013:3}
Let $\alpha$ be a number such that $\alpha/\pi$ is not a rational number.
\begin{enumerate}
    \item Show that for any integer $k$:
    \[
    \lim_{N\to\infty} \frac{1}{N} \sum_{n=1}^N e^{ik(x+n\alpha)} = \frac{1}{2\pi}\int_{-\pi}^{\pi} e^{ikt}\,dt.
    \]
    \item For every continuous periodic function $f:\mathbb{R}\to\mathbb{C}$ of period $2\pi$, we have
    \[
    \lim_{N\to\infty} \frac{1}{N}\sum_{n=1}^N f(x+n\alpha) = \frac{1}{2\pi}\int_{-\pi}^{\pi} f(t)\,dt.
    \]
\end{enumerate}
\end{exercise}

\begin{solution}
\textbf{Prerequisites:}
\begin{itemize}
    \item Geometric Series.
    \item Stone-Weierstrass Theorem (Trigonometric approximation).
    \item Weyl's Criterion for equidistribution.
\end{itemize}

\textbf{Steps:}
\begin{enumerate}
    \item \textbf{Case $k=0$:} The left side is $\frac{1}{N} \sum 1 = 1$. The right side is $\frac{1}{2\pi} \int_{-\pi}^\pi 1 \, dt = 1$. The equality holds.\\[2ex]
    \item \textbf{Case $k \neq 0$:} The right side integral is $\frac{1}{2\pi} \left[ \frac{e^{ikt}}{ik} \right]_{-\pi}^\pi = 0$ since $e^{ik\pi} = e^{-ik\pi} = (-1)^k$.
    The left side is a geometric series:
    \[
    \frac{e^{ikx}}{N} \sum_{n=1}^N (e^{ik\alpha})^n = \frac{e^{ikx}}{N} e^{ik\alpha} \frac{1 - e^{ikN\alpha}}{1 - e^{ik\alpha}}.
    \]\\[2ex]
    \item Since $\alpha/\pi \notin \mathbb{Q}$, $k\alpha$ is not a multiple of $2\pi$, so $e^{ik\alpha} \neq 1$. The term $\frac{e^{ikx} e^{ik\alpha} (1 - e^{ikN\alpha})}{1 - e^{ik\alpha}}$ is bounded. Thus, its limit as $N \to \infty$ is 0.\\[2ex]
    \item \textbf{General function $f$:} By Stone-Weierstrass, any continuous periodic $f$ can be uniformly approximated by trigonometric polynomials $P_\epsilon(t) = \sum c_k e^{ikt}$.\\[2ex]
    \item Write the difference as:
    \[
    \left| \frac{1}{N} \sum f - \int f \right| \leq \left| \frac{1}{N} \sum (f - P_\epsilon) \right| + \left| \frac{1}{N} \sum P_\epsilon - \int P_\epsilon \right| + \left| \int (P_\epsilon - f) \right|.
    \]\\[2ex]
    \item The first and third terms are small by uniform approximation ($\epsilon$). The middle term vanishes as $N \to \infty$ by step 2.
\end{enumerate}

\textbf{Intuition:}
This result is about "averaging". If you jump around a circle by an irrational angle $\alpha$, your positions $n\alpha$ will eventually "fill up" the circle evenly. The average of a function over these sampled points will then converge to the average of the function over the entire circle.
\end{solution}

\begin{exercise}
\label{analysis:2013:4}
Let $u$ be a positive harmonic function over the punctured complex plane $\mathbb{C}\setminus\{0\}$. Show that $u$ must be a constant function.
\end{exercise}

\begin{solution}
\textbf{Prerequisites:}
\begin{itemize}
    \item Harmonic functions: Liouville's Theorem (a bounded-below harmonic function on $\mathbb{R}^n$ is constant).
    \item Complex Analysis: Universal covering maps (e.g., $e^z : \mathbb{C} \to \mathbb{C} \setminus \{0\}$).
\end{itemize}

\textbf{Steps:}
\begin{enumerate}
    \item Consider the universal covering map $p: \mathbb{C} \to \mathbb{C} \setminus \{0\}$ defined by $p(z) = e^z$.
    \item Lift the function $u$ to the covering space by defining $v(z) = u(e^z)$.
    \item Since $u$ is harmonic on $\mathbb{C} \setminus \{0\}$ and $e^z$ is holomorphic, the composition $v = u \circ e^z$ is harmonic on the entire complex plane $\mathbb{C} \cong \mathbb{R}^2$.
    \item We are given that $u > 0$, which implies $v > 0$ for all $z \in \mathbb{C}$.
    \item By Liouville's Theorem for harmonic functions, any harmonic function on $\mathbb{R}^2$ that is bounded from below must be constant.
    \item Since $v$ is constant and the map $e^z$ is surjective onto $\mathbb{C} \setminus \{0\}$, it follows that $u$ must also be constant.
\end{enumerate}

\textbf{Intuition:}
A positive harmonic function on the whole plane is like a "balanced" surface that never dips below zero; the only way to stay balanced and bounded is to be flat (constant). By using the exponential map, we "unroll" the punctured plane into the whole plane, allowing us to use this powerful property.
\end{solution}

\begin{exercise}
\label{analysis:2013:5}
Suppose $H=L^2(B)$, $B$ is the unit ball in $\mathbb{R}^d$. Let $K(x,y)$ be a measurable function on $B\times B$ that satisfies
\[
|K(x,y)| \leq A|x-y|^{-d+\alpha}
\]
for some $\alpha>0$, whenever $x,y\in B$. Define
\[
Tf(x) = \int_B K(x,y)f(y)\,dy.
\]
\begin{enumerate}
    \item Prove that $T$ is a bounded operator on $H$.
    \item Prove that $T$ is compact.
\end{enumerate}
\end{exercise}

\begin{solution}
\textbf{Prerequisites:}
\begin{itemize}
    \item Functional Analysis: Schur's Test for boundedness.
    \item Compact Operators: Limits of Hilbert-Schmidt operators.
    \item Integral Calculus: Integration in polar coordinates in $\mathbb{R}^d$.
\end{itemize}

\textbf{Steps:}
\begin{enumerate}
    \item \textbf{Boundedness via Schur's Test:} An operator $T$ with kernel $K$ is bounded on $L^2$ if $\int_B |K(x,y)| \, dy$ and $\int_B |K(x,y)| \, dx$ are bounded by a constant $C$.

    \item Compute the integral:
    \[
    \int_B |x-y|^{-d+\alpha} \, dy \leq \int_{|z| \leq 2} |z|^{-d+\alpha} \, dz = \omega_{d-1} \int_0^2 r^{-d+\alpha} r^{d-1} \, dr = \omega_{d-1} \frac{2^\alpha}{\alpha} < \infty.
    \]
    Thus, $T$ is bounded.

    \item \textbf{Compactness:} Define $K_n(x,y) = K(x,y)$ if $|x-y| \geq 1/n$ and $0$ otherwise. Since $K_n$ is bounded on a set of finite measure, it is in $L^2(B \times B)$, so $T_n$ is Hilbert-Schmidt and thus compact.

    \item Show $T_n \to T$ in operator norm. The difference $T - T_n$ has kernel $K$ restricted to $|x-y| < 1/n$. By Schur's Test:
    \[
    \|T - T_n\| \leq \int_{|z| < 1/n} A |z|^{-d+\alpha} \, dz = \frac{A \omega_{d-1}}{\alpha n^\alpha} \to 0 \text{ as } n \to \infty.
    \]

    \item Since $T$ is the limit of compact operators, $T$ is compact.
\end{enumerate}

\textbf{Intuition:}
The kernel $K(x,y)$ has a singularity at $x=y$, but it's "weak" enough (the power $-d+\alpha$ is better than $-d$) that it's still integrable. This integrability ensures the operator is well-behaved (bounded). By chopping off the singularity, we get operators that are very "smooth" (Hilbert-Schmidt/Compact), and the whole operator is just a limit of these smooth ones.
\end{solution}

\begin{exercise}
\label{analysis:2013:6}
Let $A$ be an $n\times n$ real non-degenerate symmetric matrix. For $\lambda\in\mathbb{R}^+$, we define:
\[
\int_{-\infty}^{\infty} \exp\left(i\lambda x^2\right)dx = \lim_{\epsilon\to 0^+}\int_{-\infty}^{\infty} \exp\left(i\lambda x^2 - \frac{1}{2}\epsilon x^2\right)dx.
\]
Show that:
\[
\int_{\mathbb{R}^n} \exp\left(i\frac{\lambda}{2}\langle Ax,x\rangle - i\langle x,\xi\rangle\right)dx = \left(\frac{2\pi}{\lambda}\right)^{n/2}|\det A|^{-1/2}\exp\left(-\frac{i}{2\lambda}\langle A^{-1}\xi,\xi\rangle\right)\exp\left(\frac{i\pi}{4}\operatorname{sgn}A\right).
\]
\end{exercise}

\begin{solution}
\textbf{Prerequisites:}
\begin{itemize}
    \item Multivariable Calculus: Change of variables and diagonalization of symmetric matrices.
    \item Fresnel Integrals: $\int_{-\infty}^\infty e^{iax^2} dx = \sqrt{\frac{\pi}{|a|}} e^{i\frac{\pi}{4}\text{sgn}(a)}$.
    \item Fourier Transform of a Gaussian.
\end{itemize}

\textbf{Steps:}
\begin{enumerate}
    \item Diagonalize $A = Q D Q^T$, where $D = \text{diag}(\mu_1, \dots, \mu_n)$ are the eigenvalues. Change variables $x = Qy$, so $dx = dy$.

    \item The exponent becomes $i\frac{\lambda}{2} \sum \mu_j y_j^2 - i \sum \eta_j y_j$, where $\eta = Q^T \xi$. The integral factorizes into $n$ 1D integrals:
    \[
    \prod_{j=1}^n \int_{-\infty}^\infty \exp\left( i \frac{\lambda \mu_j}{2} y_j^2 - i \eta_j y_j \right) dy_j.
    \]

    \item Complete the square: $i \frac{\lambda \mu_j}{2} y_j^2 - i \eta_j y_j = i \frac{\lambda \mu_j}{2} (y_j - \frac{\eta_j}{\lambda \mu_j})^2 - \frac{i \eta_j^2}{2 \lambda \mu_j}$.

    \item Use the 1D result: $\int_{-\infty}^\infty \exp( i \frac{\lambda \mu_j}{2} z^2 ) dz = \sqrt{\frac{2\pi}{\lambda |\mu_j|}} e^{i \frac{\pi}{4} \text{sgn}(\mu_j)}$.

    \item Multiply the results: The product of $\sqrt{1/|\mu_j|}$ gives $|\det A|^{-1/2}$, the sum of $\text{sgn}(\mu_j)$ gives $\text{sgn}(A)$, and the sum $\sum \frac{\eta_j^2}{\mu_j}$ gives $\langle A^{-1} \xi, \xi \rangle$.
\end{enumerate}

\textbf{Intuition:}
This is a multidimensional version of the Gaussian integral. Just as a 1D Gaussian is determined by its variance, this integral is determined by the "variances" along the principal axes of the matrix $A$. The phase shift (the $\operatorname{sgn}A$ term) comes from whether each axis corresponds to a "valley" or a "peak" in the complex oscillation.
\end{solution}

\subsection{Team}

\begin{exercise}
\label{analysis:2013:7}
Suppose $\Delta = \{z \in \mathbb{C} \mid |z| < 1\}$ is the open unit disk in the complex plane. Show that for any holomorphic function $f: \Delta \to \Delta$,
\[
\frac{|f'(z)|}{1 - |f(z)|^2} \leq \frac{1}{1 - |z|^2} \tag{1}
\]
for all $z$ in $\Delta$. If equality holds in (1) for some $z_0 \in \Delta$, show that $f \in \operatorname{Aut}(\Delta)$, and that the equality holds for all $z \in \Delta$.
\end{exercise}

\begin{solution}
\textbf{Prerequisites:}
\begin{itemize}
    \item Complex Analysis: Schwarz Lemma.
    \item Möbius Transformations: Automorphisms of the unit disk $\Delta$.
    \item Chain Rule for holomorphic functions.
\end{itemize}

\textbf{Steps:}
\begin{enumerate}
    \item Let $\psi_a(z) = \frac{a-z}{1-\bar{a}z}$ be an automorphism of $\Delta$ that maps $a$ to $0$.

    \item For a fixed $z_0 \in \Delta$, define a new function $g(w) = \psi_{f(z_0)} \circ f \circ \psi_{z_0}^{-1}(w)$.

    \item Note that $g(0) = \psi_{f(z_0)}(f(z_0)) = 0$, and since $f: \Delta \to \Delta$, we have $g: \Delta \to \Delta$.

    \item By the Schwarz Lemma, $|g'(0)| \leq 1$. We calculate $g'(0)$ using the chain rule:
    \[
    |g'(0)| = |\psi_{f(z_0)}'(f(z_0))| \cdot |f'(z_0)| \cdot |(\psi_{z_0}^{-1})'(0)| = \frac{1}{1-|f(z_0)|^2} \cdot |f'(z_0)| \cdot (1-|z_0|^2).
    \]

    \item The inequality $|g'(0)| \leq 1$ yields $\frac{|f'(z_0)|}{1-|f(z_0)|^2} \leq \frac{1}{1-|z_0|^2}$.

    \item If equality holds at $z_0$, then $|g'(0)| = 1$. By the Schwarz Lemma, $g(w) = e^{i\theta}w$, which is an automorphism. Thus $f$ is a composition of automorphisms, hence $f \in \operatorname{Aut}(\Delta)$.
\end{enumerate}

\textbf{Intuition:}
The Schwarz-Pick theorem says that holomorphic maps from the disk to itself are "contractions" with respect to the hyperbolic metric. The expression $\frac{|f'(z)|}{1-|f(z)|^2} dz$ is the hyperbolic length of the image of a small vector at $z$. The theorem states that this image vector is never longer than the original vector in the hyperbolic sense.
\end{solution}

\begin{exercise}
\label{analysis:2013:8}
Let $f$ be a function of bounded variation on $[a, b]$, $f_1$ its generalized derivative as a measure, i.e. $f(x) - f(a) = \int_a^x f_1(y) dy$ for every $x \in [a, b]$ and $f_1(x)$ is an integrable function on $[a, b]$. Let $f'$ be its weak derivative as a generalized function, i.e. $\int_a^b f(x) g'(x) dx = -\int_a^b f'(x) g(x) dx$, for any smooth function $g(x)$ on $[a, b], g(a) = g(b) = 0$. Show that:
\begin{enumerate}
    \item If $f$ is absolutely continuous, then $f' = f_1$.
    \item If the weak derivative $f'$ of $f$ is an integrable function on $[a, b]$, then $f(x)$ is equal to an absolutely continuous function outside a set of measure zero.
\end{enumerate}
\end{exercise}

\begin{solution}
\textbf{Prerequisites:}
\begin{itemize}
    \item Real Analysis: Absolutely continuous functions and the Fundamental Theorem of Calculus for Lebesgue integrals.
    \item Theory of Distributions: Definition of weak derivatives.
    \item Integration by Parts for Lebesgue integrals.
\end{itemize}

\textbf{Steps:}
\begin{enumerate}
    \item \textbf{Part (a):} If $f$ is absolutely continuous, then $f(x) = f(a) + \int_a^x f_1(t) dt$. Its pointwise derivative $f_1$ exists almost everywhere and is integrable.

    \item Using integration by parts for absolutely continuous functions:
    \[
    \int_a^b f(x) g'(x) dx = [f(x)g(x)]_a^b - \int_a^b f_1(x) g(x) dx = -\int_a^b f_1(x) g(x) dx.
    \]
    Since this holds for all test functions $g$, the weak derivative $f'$ coincides with $f_1$.

    \item \textbf{Part (b):} Suppose the weak derivative $f'$ is integrable. Define an absolutely continuous function $F(x) = \int_a^x f'(t) dt$.

    \item By the result in (a), the weak derivative of $F$ is $f'$. Thus, for any test function $g$:
    \[
    \int_a^b (f(x) - F(x)) g'(x) dx = \int_a^b f g' - \int_a^b F g' = -\int_a^b f' g + \int_a^b f' g = 0.
    \]

    \item A distribution whose derivative is zero must be a constant. Thus $f(x) - F(x) = C$ almost everywhere. This means $f(x) = F(x) + C$ a.e., which is absolutely continuous.
\end{enumerate}

\textbf{Intuition:}
Weak derivatives generalize the concept of a derivative to functions that might not be differentiable everywhere. However, if the weak derivative happens to be a nice integrable function, then the original function must have been "nice" too—specifically, absolutely continuous, which is the strongest form of continuity allowing for the Fundamental Theorem of Calculus.
\end{solution}

\begin{exercise}
\label{analysis:2013:9}
Show that the convex hull of the roots of any polynomial contains all its critical points as well as all the zeros of higher derivatives of the polynomial.
\end{exercise}

\begin{solution}
\textbf{Prerequisites:}
\begin{itemize}
    \item Algebra: Logarithmic derivative of a polynomial.
    \item Geometry: Definition of convex hull and convex combinations.
    \item Complex Numbers: Interpretation of $\sum w_i (z - z_i) = 0$ as $z$ being a center of mass.
\end{itemize}

\textbf{Steps:}
\begin{enumerate}
    \item Let $P(z) = a \prod_{i=1}^n (z - z_i)$. Its logarithmic derivative is $\frac{P'(z)}{P(z)} = \sum_{i=1}^n \frac{1}{z - z_i}$.

    \item If $z$ is a root of $P'(z)$, then either $P(z) = 0$ (so $z$ is one of the $z_i$, which is in the convex hull), or $\sum_{i=1}^n \frac{1}{z - z_i} = 0$.

    \item Take the complex conjugate: $\sum_{i=1}^n \frac{1}{\bar{z} - \bar{z}_i} = \sum_{i=1}^n \frac{z - z_i}{|z - z_i|^2} = 0$.

    \item Rearrange the terms: $z \sum_{i=1}^n \frac{1}{|z - z_i|^2} = \sum_{i=1}^n \frac{z_i}{|z - z_i|^2}$.

    \item This shows $z = \sum w_i z_i$ where $w_i = \frac{1/|z-z_i|^2}{\sum 1/|z-z_k|^2}$ are positive weights that sum to 1. Thus $z$ is in the convex hull of $\{z_i\}$.

    \item By induction, since the roots of $P^{(k)}$ are in the convex hull of the roots of $P^{(k-1)}$, they are all contained in the convex hull of the original roots.
\end{enumerate}

\textbf{Intuition:}
This is the Gauss-Lucas Theorem. It tells us that the "descendants" of a polynomial (its derivatives) are always "trapped" within the geometric boundaries set by their "ancestors" (the original roots). Physically, you can think of the roots as point charges; the critical points are the equilibrium positions where the total force vanishes, which must lie within the region bounded by the charges.
\end{solution}

\begin{exercise}
\label{analysis:2013:10}
Let $D \subset \mathbb{R}^3$ be an open domain. Show that every smooth vector field $\mathbf{F} = (P, Q, R)$ over $D$ can be written as $\mathbf{F} = \mathbf{F}_1 + \mathbf{F}_2$ such that $\operatorname{rot}(\mathbf{F}_1) = 0$, $\operatorname{div}(\mathbf{F}_2) = 0$.
\end{exercise}

\begin{solution}
\textbf{Prerequisites:}
\begin{itemize}
    \item Vector Calculus: Helmholtz Decomposition.
    \item Differential Equations: Poisson's equation $\Delta \phi = \rho$.
    \item Integral Calculus: Newton potential formula.
\end{itemize}

\textbf{Steps:}
\begin{enumerate}
    \item We seek a decomposition $\mathbf{F} = \nabla \phi + \nabla \times \mathbf{A}$. Set $\mathbf{F}_1 = \nabla \phi$ and $\mathbf{F}_2 = \nabla \times \mathbf{A}$.

    \item To satisfy $\operatorname{div} \mathbf{F}_2 = 0$, we only need $\mathbf{F}_2$ to be a curl. To satisfy $\operatorname{rot} \mathbf{F}_1 = 0$, we need $\mathbf{F}_1$ to be a gradient.

    \item Taking the divergence of $\mathbf{F} = \mathbf{F}_1 + \mathbf{F}_2$, we get $\operatorname{div} \mathbf{F} = \operatorname{div}(\nabla \phi) = \Delta \phi$.

    \item We solve the Poisson equation $\Delta \phi = \operatorname{div} \mathbf{F}$. In $\mathbb{R}^3$, the solution is given by the Newton potential:
    \[
    \phi(x) = -\frac{1}{4\pi} \int_{\mathbb{R}^3} \frac{\operatorname{div} \mathbf{F}(y)}{|x-y|} dy.
    \]

    \item Once $\phi$ is found, we define $\mathbf{F}_1 = \nabla \phi$. Then define $\mathbf{F}_2 = \mathbf{F} - \mathbf{F}_1$.

    \item By construction, $\operatorname{div} \mathbf{F}_2 = \operatorname{div} \mathbf{F} - \Delta \phi = 0$, and $\operatorname{rot} \mathbf{F}_1 = \operatorname{rot} \nabla \phi = 0$.
\end{enumerate}

\textbf{Intuition:}
The Helmholtz decomposition states that any sufficiently smooth vector field can be split into an "irrotational" (curl-free) part and a "solenoidal" (divergence-free) part. Think of the field as a fluid flow: $\mathbf{F}_1$ represents the part of the flow caused by sources and sinks, while $\mathbf{F}_2$ represents the part caused by swirling vortices.
\end{solution}

\begin{exercise}
\label{analysis:2013:11}
Let $\mathbf{H}$ be a Hilbert space and $\mathbf{A}$ a compact self-adjoint linear operator over $\mathbf{H}$. Show that there exists an orthonormal basis of $\mathbf{H}$ consisting of eigenvectors $\varphi_n$ of $\mathbf{A}$ with non-zero eigenvalues $\lambda_n$ such that every vector $\xi \in \mathbf{H}$ can be written as: $\xi = \sum_k c_k \varphi_k + \xi'$, where $\xi' \in \operatorname{Ker} \mathbf{A}$. We also have $\mathbf{A}\xi = \sum_k \lambda_k c_k \varphi_k$. If there are infinitely many eigenvectors, then $\lim_{n \to \infty} \lambda_n = 0$.
\end{exercise}

\begin{solution}
\textbf{Proof, with the compact spectral theorem reduced to its key lemma.}

\textbf{Lemma.}
If \(A\) is a nonzero compact self-adjoint operator on a Hilbert space \(H\), then \(A\) has a nonzero eigenvalue.

\textit{Proof of the lemma.}
For a compact self-adjoint operator, one first shows that the numerical radius is attained on the unit sphere. Indeed, choose unit vectors \(x_j\) such that
\[
|\langle Ax_j,x_j\rangle|\to
\sup_{\|x\|=1}|\langle Ax,x\rangle|.
\]
After passing to a weakly convergent subsequence \(x_j\rightharpoonup x\), compactness gives \(Ax_j\to Ax\) in norm. The standard polarization identity then shows that equality in the extremal variational problem forces
\[
Ax=\lambda x
\]
for a real number \(\lambda\) with \(|\lambda|=\sup_{\|y\|=1}|\langle Ay,y\rangle|\). Since \(A\neq0\), this \(\lambda\neq0\). This proves the lemma.

Now construct the decomposition. Pick a nonzero eigenvalue \(\lambda_1\) and its eigenspace \(E_{\lambda_1}\). If \(\lambda\neq\mu\), then for \(u\in E_\lambda\) and \(v\in E_\mu\),
\[
\lambda\langle u,v\rangle
=\langle Au,v\rangle
=\langle u,Av\rangle
=\mu\langle u,v\rangle,
\]
so \(E_\lambda\perp E_\mu\). Also \(E_\lambda\) is finite-dimensional when \(\lambda\neq0\): if it were infinite-dimensional, its unit ball would contain an orthonormal sequence \(\{e_j\}\), and \(Ae_j=\lambda e_j\) would have no convergent subsequence, contradicting compactness.

Remove \(E_{\lambda_1}\) and repeat the argument on the orthogonal complement. The complement is invariant under \(A\), because \(A\) is self-adjoint. This gives a countable orthonormal family of eigenvectors for all nonzero eigenspaces. Let \(M\) be the closed span of those eigenspaces.

The restriction of \(A\) to \(M^\perp\) is still compact and self-adjoint. It has no nonzero eigenvalue by construction. By the lemma, this restriction must be the zero operator. Hence
\[
M^\perp\subset \ker A.
\]
Therefore every \(\xi\in H\) decomposes uniquely as
\[
\xi=\sum_k c_k\varphi_k+\xi',
\qquad
\xi'\in\ker A,
\]
where \(c_k=\langle \xi,\varphi_k\rangle\), and the series converges in the Hilbert norm. Applying \(A\) term by term gives
\[
A\xi=\sum_k \lambda_k c_k\varphi_k.
\]

Finally, if infinitely many nonzero eigenvalues did not tend to \(0\), then there would be \(\epsilon>0\) and an orthonormal sequence \(\varphi_{k_j}\) with \(|\lambda_{k_j}|\geq\epsilon\). The sequence
\[
A\varphi_{k_j}=\lambda_{k_j}\varphi_{k_j}
\]
has no norm-convergent subsequence, contradicting compactness. Hence \(\lambda_k\to0\) along any infinite list of nonzero eigenvalues.
\end{solution}

\begin{exercise}
\label{analysis:2013:12}
A function $f : \mathbb{R} \to \mathbb{R}$ is called convex if $f(\lambda x + (1 - \lambda) x') \leq \lambda f(x) + (1 - \lambda) f(x')$ for $0 \leq \lambda \leq 1$ and each $x, x' \in \mathbb{R}$. We assume $|f(x)| < \infty$ for all $x$.
\begin{enumerate}
    \item Show that a convex function $f$ is continuous and the function $g(y) = \max_{x \in \mathbb{R}} (xy - f(x))$ is a well-defined convex function over $\mathbb{R}$.
    \item Show that a convex function $f$ is differentiable except at most countably many points.
    \item $f$ is differentiable everywhere if both $f$ and $g$ are strictly convex.
\end{enumerate}
\end{exercise}

\begin{solution}
\textbf{As written, part (a) is missing a growth hypothesis.}
The continuity assertion is correct, but the claim that
\[
g(y)=\max_{x\in\mathbb R}(xy-f(x))
\]
is well-defined for every finite convex \(f\) is false. For example, if \(f(x)\equiv0\), then \(xy-f(x)=xy\), whose maximum over \(x\in\mathbb R\) is not finite unless \(y=0\), and even at \(y=0\) the maximum is attained by every \(x\). Thus the intended transform needs an additional coercivity condition, or the word ``max'' should be replaced by the extended-valued supremum.

\textbf{Corrected Legendre-transform version.}
If \(f\) is finite convex and superlinear, meaning
\[
\frac{f(x)}{|x|}\to+\infty
\qquad \text{as } |x|\to\infty,
\]
then \(xy-f(x)\to-\infty\) as \(|x|\to\infty\) for each fixed \(y\). Since \(f\) is continuous, the maximum is attained and finite. Moreover \(g\) is convex, because it is the supremum of affine functions of \(y\):
\[
g(\lambda y_1+(1-\lambda)y_2)
\leq \lambda g(y_1)+(1-\lambda)g(y_2).
\]

\textbf{Continuity of convex functions.}
For any compact interval \([a,b]\), the secant slopes of a finite convex function are monotone. If \(a<x<y<z<b\), then
\[
\frac{f(y)-f(x)}{y-x}
\leq
\frac{f(z)-f(x)}{z-x}
\leq
\frac{f(z)-f(a)}{z-a},
\]
and a similar lower bound holds from the left. Thus the one-sided slopes are locally bounded, so \(f\) is locally Lipschitz on compact subintervals and hence continuous.

\textbf{Differentiability except countably many points.}
For each \(x\), the one-sided derivatives
\[
f'_-(x)\leq f'_+(x)
\]
exist as extended monotone limits, and both are nondecreasing functions of \(x\). The function \(f\) is differentiable exactly when \(f'_-(x)=f'_+(x)\). If \(f'_-(x)<f'_+(x)\), choose a rational number \(q_x\) in the interval \((f'_-(x),f'_+(x))\). Monotonicity of the one-sided derivatives makes the intervals \((f'_-(x),f'_+(x))\) disjoint for distinct nondifferentiability points. Hence the map \(x\mapsto q_x\) is injective into \(\mathbb Q\), so there are at most countably many such points.

\textbf{Strict convexity of the dual forces differentiability.}
Assume the corrected Legendre transform is finite and that both \(f\) and \(g\) are strictly convex. If \(f\) were not differentiable at \(x_0\), then its subdifferential would contain a nontrivial interval \([a,b]\) with \(a<b\). For every \(y\in[a,b]\), the line of slope \(y\) supports \(f\) at \(x_0\), so
\[
g(y)=x_0y-f(x_0)
\]
on \([a,b]\). This is affine on a nontrivial interval, contradicting strict convexity of \(g\). Therefore \(f\) is differentiable everywhere.
\end{solution}
